\documentclass{training}
\schoollogo{../../../resources/latex/crack.pdf}
\usepackage[utf8]{inputenc}
\usepackage[T1]{fontenc}
\usepackage{lmodern}
\setlang{french}
\usepackage{hyperref}
\usepackage{graphicx}
\usepackage{tabularx}
\usepackage{fvextra}
\fvset{breaklines=true,breakanywhere=true,breaksymbol={},breaksymbolleft={},breaksymbolright={},breakanywheresymbolpre={},breakanywheresymbolpost={}}
\makeatletter
\g@addto@macro\UrlBreaks{%
  \do\a\do\b\do\c\do\d\do\e\do\f\do\g\do\h\do\i\do\j\do\k\do\l\do\m%
  \do\n\do\o\do\p\do\q\do\r\do\s\do\t\do\u\do\v\do\w\do\x\do\y\do\z%
  \do\A\do\B\do\C\do\D\do\E\do\F\do\G\do\H\do\I\do\J\do\K\do\L\do\M%
  \do\N\do\O\do\P\do\Q\do\R\do\S\do\T\do\U\do\V\do\W\do\X\do\Y\do\Z%
  \do\0\do\1\do\2\do\3\do\4\do\5\do\6\do\7\do\8\do\9%
  \do\<\do\>\do\(\do\)\do\'\do\/}
\schoollogo{../../../resources/latex/crack.pdf}
\makeatother


\begin{document}

\title{Warm-ups croisés : Node.js, MINA, FreeBSD --- A-EPI-006}
\subtitle{Trois domaines, trois familles, aucune excuse}
\member{Ted ``Tedsig42'' Kouhouenou}{ted.sig42@tutamail.com}
\summary{Module de raccord : couvrir les exemples du papier qui ne sont pas des CVE du noyau Linux. Node.js path.normalize, Apache MINA ObjectSerializationDecoder et les mappings FreeBSD/XNU montrent que les trois familles ne dépendent ni d'un langage, ni d'un éditeur, ni d'un noyau. Aucun exploit : uniquement des prédicats, des machines à états et des matrices de branches.}
\version{1.0}
\prerequisite{M-EPI-001, M-EPI-002, bases de JavaScript et de Java}
\maketitle

\newpage
\tableofcontents
\newpage

\chapter{Instructions}
\begin{itemize}
  \item Deux jours.
  \item Aucune cible réelle, aucun payload de désérialisation, aucun chemin
        traversé sur un vrai système de fichiers.
  \item Les tests Node.js se font avec \texttt{path.win32} sur des chaînes,
        sans appel à \texttt{fs}.
  \item Pour MINA, lisez les avis, les sources et les diffs ; ne construisez
        pas de flux sérialisé malveillant.
  \item Pour FreeBSD et XNU, travaillez à partir des avis publics et de la
        section 18.11 du papier. Distinguez ce qui est démontré de ce qui
        est seulement cartographié.
\end{itemize}

\subsection*{Format du rendu --- en-tête obligatoire}

Tout fichier rendu dans ce module doit commencer par l'en-tête
ci-dessous, \textbf{adapté à la syntaxe de commentaire du langage}
utilisé.

\subsubsection*{Modèle générique}

\begin{minted}{c}
/**
 * Module: <code module>
 * Exercise: <numéro exo>
 * File: <nom du fichier>
 * Author: <votre login JeuneCrack>
 *
 * Description: <une phrase, pas de roman>
 *
 * Return: <voir ci-dessous>
 */
\end{minted}

\subsubsection*{Login}

Votre identifiant \textbf{JeuneCrack} (ex: \texttt{tedsig42}).
Pour un rendu en binôme : \texttt{login1/login2}.
Ajoutez un handle Matrix après le login si vous le souhaitez
(ex: \texttt{tedsig42 @tedsig42}).

\subsubsection*{Return --- selon le type de fichier}

\begin{tabularx}{\linewidth}{lX}
  \toprule
  \textbf{Type de fichier} & \textbf{Valeur de Return} \\
  \midrule
  C (\texttt{.c}) &
  \texttt{return (0);} (succès),
  \texttt{return (-1);} (erreur),
  \texttt{return (valeur);} (résultat calculé) \\

  Python (\texttt{.py}) &
  \texttt{0} (succès),
  \texttt{1} (erreur),
  \texttt{N/A} (analyse, script sans exit) \\

  Shell (\texttt{.sh}) &
  \texttt{exit 0} (succès),
  \texttt{exit 1} (erreur) \\

  Markdown (\texttt{.md}) &
  \texttt{N/A} --- les fichiers d'analyse n'ont pas de valeur de retour \\

  Exploit C/Python &
  \texttt{shell} (shell obtenu),
  \texttt{flag: <flag>} (flag capturé),
  \texttt{proof: <hash>} (preuve de concept) \\

  Binaire (\texttt{.bin}) &
  \texttt{N/A} --- pas d'en-tête pour les fichiers binaires \\
  \bottomrule
\end{tabularx}

\subsubsection*{Exemples par langage}

\textbf{C (\texttt{.c})} :
\begin{minted}{c}
/**
 * Module: F-EXP-001
 * Exercise: ex02
 * File: exploit.c
 * Author: tedsig42
 *
 * Description: Buffer overflow with ret2libc
 *
 * Return: 0
 */
\end{minted}

\textbf{Python (\texttt{.py})} :
\begin{minted}{python}
# Module: F-WEB-001
# Exercise: ex05
# File: ssrf-exploit.py
# Author: tedsig42
#
# Description: Blind SSRF bypass via redirect
#
# Return: 0
\end{minted}

\textbf{Shell (\texttt{.sh})} :
\begin{minted}{bash}
# Module: F-SUP-001
# Exercise: ex03
# File: dependency-confusion.sh
# Author: tedsig42
#
# Description: Automated dependency confusion test
#
# Return: exit 0
\end{minted}

\textbf{Markdown (\texttt{.md})} --- analyse, rapport :
\begin{Verbatim}[fontsize=\normalsize]
<!--
Module: M-SRC-001
Exercise: ex04
File: audit-report.md
Author: tedsig42

Description: Code audit of vulnerable-httpd

Return: N/A
-->
\end{Verbatim}

\chapter{Instructions Communes}

\section{Structure du D\'ep\^ot}

Votre d\'ep\^ot Git doit respecter la convention de nommage suivante pour
tous les modules soumis. L'outil d'\'evaluation (\texttt{pendragon})
se base sur cette structure pour localiser et tester votre travail.

\begin{console}
<login>/
+-- <MODULE-ID>-<slug>/         # ex. M-BLS-001-business-logic/
|   +-- <turnindir>/            # ex. ex00-pricing-manipulation/
|   |   +-- <file1>             # ex. poc-pricing.py
|   |   +-- <file2>             # ex. report.md
|   |   +-- ...
|   +-- <turnindir>/            # ex. ex01-race-conditions/
|   |   +-- ...
|   +-- ...
+-- ...
\end{console}

\begin{itemize}
  \item \textbf{R\'epertoire du module} : le nom complet du r\'epertoire tel
        qu'affich\'e dans votre curriculum (ex. \texttt{M-BLS-001-business-logic}).
  \item \textbf{R\'epertoire de rendu} : la valeur de \texttt{\textbackslash turnindir}
        fournie dans chaque exercice (ex. \texttt{ex00-pricing-manipulation}).
  \item \textbf{Fichiers de l'exercice} : exactement les fichiers list\'es dans
        \texttt{\textbackslash exfiles}. Ni fichiers suppl\'ementaires, ni fichiers manquants.
\end{itemize}

\section{Qu'est-ce qu'une Preuve de Concept Valide}

Une PoC doit \^etre :

\begin{enumerate}
  \item \textbf{Reproductible}. N'importe qui avec le m\^eme environnement peut
        l'ex\'ecuter et observer la vuln\'erabilit\'e. Les chemins cod\'es en dur,
        les adresses IP absolues et les jetons expir\'es ne sont pas acceptables.
  \item \textbf{Document\'ee}. Chaque \'etape est expliqu\'ee. L'\'evaluteur ne
        devrait pas avoir \`a deviner ce que fait le script.
  \item \textbf{Cibl\'ee}. L'exercice sp\'ecifie une cible autoris\'ee
        (lab, maquette, binaire local ou plateforme d\'esign\'ee). Ne testez
        jamais les PoC sur des syst\`emes de production qui ne vous appartiennent pas.
  \item \textbf{Minimale}. La PoC d\'emontre \emph{uniquement} la vuln\'erabilit\'e
        en question. Les charges utiles non li\'ees et le bruit sont p\'enalis\'es.
\end{enumerate}

\begin{warning}
  Une PoC qui modifie silencieusement l'\'etat du syst\`eme, supprime des
  fichiers ou persiste des changements en dehors du p\'erim\`etre de l'exercice
  sera rejet\'ee. Si votre exploit n\'ecessite un nettoyage, documentez les
  \'etapes de nettoyage dans votre rapport.
\end{warning}

\section{Processus de Soumission}

\begin{enumerate}
  \item R\'ealisez les livrables de l'exercice sur votre machine locale.
  \item Committez et poussez vers votre d\'ep\^ot Git assign\'e.
  \item Le d\'ep\^ot est gel\'e apr\`es la date limite.
  \item Les pairs clonent votre d\'ep\^ot et \'evaluent en utilisant :
        \begin{itemize}
          \item \texttt{pendragon} pour l'ex\'ecution automatis\'ee des tests
                (l\`a o\`u le module fournit des configurations mouli).
          \item La grille manuelle pour la qualit\'e du rapport, le style de
                code et la m\'ethodologie.
        \end{itemize}
\end{enumerate}

\hint{Ex\'ecutez toujours \texttt{pendragon run <MODULE-ID> <ex\_id>
<path/to/student/repo>} avant votre commit final pour v\'erifier
que votre soumission passe les v\'erifications automatis\'ees.}

\subsection*{Test avec pendragon}

Si le module inclut une configuration mouli (fichiers JSON dans
\texttt{resources/mouli/}), vous pouvez tester votre soumission localement :

\begin{console}
pendragon run M-BLS-001 ex00 ~/student-repos/tedsig42
\end{console}

\texttt{pendragon} r\'esout le chemin de l'exercice automatiquement en
cherchant \texttt{<mod-fullname>/<ex-dirname>/} dans le r\'epertoire
\'etudiant donn\'e.

\subsection*{V\'erification de style avec caliban}

Pour les modules qui n\'ecessitent du code C, \texttt{caliban} v\'erifie
la conformit\'e du style :

\begin{console}
caliban check ~/student-repos/tedsig42/M-BLS-001-business-logic
\end{console}

\texttt{caliban} applique la Norme 42 (taille des fonctions, nommage
des variables, constructions interdites) adapt\'ee aux contextes de
recherche de vuln\'erabilit\'es. Ex\'ecutez-le avant la soumission ---
un \'echec aux v\'erifications de style peut invalider l'exercice.

\section{R\`egle des Fichiers Suppl\'ementaires Interdits}

\begin{itemize}
  \item Votre d\'ep\^ot doit contenir \textbf{uniquement} les fichiers
        list\'es dans \texttt{\textbackslash exfiles} pour chaque exercice.
  \item Fichiers interdits : \texttt{.DS\_Store}, \texttt{\_\_pycache\_\_},
        \texttt{.venv/}, \texttt{node\_modules/}, \texttt{.git/},
        binaires compil\'es (sauf explicitement list\'es dans \texttt{exfiles}),
        fichiers d'\'echange d'\'editeur, fichiers journaux et tout fichier
        non demand\'e.
  \item Utilisez \texttt{.gitignore} pour exclure automatiquement les
        fichiers ind\'esirables.
\end{itemize}

\section{Format du Rapport}

Lorsqu'un rapport (\texttt{report.md}) est requis, il doit inclure :

\begin{enumerate}
  \item \textbf{Classe de vuln\'erabilit\'e} (identifiant CWE si applicable).
  \item \textbf{Étapes de reproduction} --- num\'erot\'ees, pr\^etes
        \`a \^etre copi\'ees-coll\'ees.
  \item \textbf{Charges utiles} --- chaînes exactes, requ\^etes ou
        entr\'ees utilis\'ees.
  \item \textbf{Impact} --- ce qu'un attaquant peut accomplir.
  \item \textbf{Mitigation} --- comment le d\'eveloppeur devrait corriger.
  \item \textbf{Sources} --- toutes les ressources consult\'ees (voir
        Section~\ref{sec:cite-sources}).
\end{enumerate}

\subsection*{En-t\^ete de Rendu}

Chaque fichier source doit inclure l'en-t\^ete de rendu d\'efini dans
\path{../../../resources/latex/deliverable-header.tex}, adapt\'e
\`a la syntaxe de commentaire du langage utilis\'e.

\section{Cibles Autoris\'ees Uniquement}

\begin{warning}
  Vous pouvez \textbf{uniquement} tester les cibles explicitement list\'ees
  dans l'exercice (labs, serveurs locaux, maquettes, binaires d\'esign\'es).
  Tester des syst\`emes de production r\'eels, des applications tierces
  ou toute cible qui ne vous appartient pas ou pour laquelle vous n'avez
  pas d'autorisation \'ecrite est interdit et peut entraîner une
  disqualification imm\'ediate.
\end{warning}

L'\'ethique du bug bounty s'applique tout au long de ce curriculum :
\begin{itemize}
  \item Aucun test sur des programmes en production sans un accord de
        p\'erim\`etre valide.
  \item Aucun d\'eni de service sauf si l'exercice l'exige explicitement.
  \item Aucune exfiltration de donn\'ees au-del\`a de ce qui est
        n\'ecessaire pour d\'emontrer la vuln\'erabilit\'e.
  \item Si vous acc\'edez accidentellement \`a des donn\'ees r\'eelles
        d'utilisateurs, arr\^etez imm\'ediatement et documentez ce qui
        s'est pass\'e.
\end{itemize}

\section{Citez Toutes Vos Sources}
\label{sec:cite-sources}

Toute ressource consult\'ee pendant votre travail doit \^etre cit\'ee.
Cela inclut :

\begin{itemize}
  \item Articles acad\'emiques et billets de blog.
  \item Rapports publics de bug bounty (HackerOne, Bugcrowd, etc.).
  \item Documentation officielle (RFC, pages de manuel, documentation
        fournisseur).
  \item \textbf{Outils d'IA} : nom du mod\`ele, r\'esum\'e de la
        question, date et m\'ethode de v\'erification (voir le chapitre
        Instructions IA pour le format de citation exact).
  \item Extraits de code adapt\'es de projets open-source (avec
        attribution de licence si requis).
\end{itemize}

\subsection*{Format de citation}

\begin{minted}{markdown}
# R\'ef\'erences

- [1] PortSwigger, "Business Logic Flaws",
  https://portswigger.net/web-security/logic-flaws
  (consult\'e le 2025-03-15).
- [2] Rapport HackerOne #123456, "IDOR on /api/users",
  anonymis\'e, consult\'e via divulgation publique.
- [3] G\'en\'er\'e avec l'assistance de Claude 3.5 Sonnet
  le 2025-03-14. Question : "Write a Python script that...".
  V\'erifi\'e par : revue manuelle de chaque requ\^ete HTTP.
\end{minted}

\subsection*{Attribution en ligne}

Si un exercice vous demande explicitement d'utiliser une technique ou
un outil sp\'ecifique, vous n'avez pas besoin de citer l'exercice
lui-m\^eme comme source. Citez tout ce qui se trouve en dehors de
l'\'enonc\'e de l'exercice.

\chapter{Instructions sur l'IA}

\section{Aperçu de la Politique}

Les outils de codage assistés par IA (ChatGPT, Claude, Copilot, Gemini,
Cursor, Codeium, et similaires) sont \textbf{autorisés} tout au long
de ce cursus. La recherche en sécurité est un domaine où la maîtrise
de l'IA devient essentielle --- apprendre à utiliser ces outils
efficacement fait partie de votre formation.

Cependant, l'objectif est la \textbf{compréhension}, pas l'automatisation.
Si vous ne pouvez pas expliquer chaque ligne de votre exploit sans
l'aide de l'IA, vous n'avez pas terminé l'exercice.

\begin{warning}
  Utiliser l'IA pour contourner l'objectif pédagogique d'un exercice
  --- par exemple, lui demander de rétro-ingéniérer un binaire pour
  ne pas avoir à apprendre le désassembleur, ou de générer une chaîne
  d'exploitation complète que vous ne comprenez pas --- va à l'encontre
  du but du cursus et sera traité comme une fraude académique.
\end{warning}

\section{Utilisations Autorisées}

\begin{itemize}
  \item \textbf{Explication de concepts}. Demander à l'IA de clarifier
        une classe de vulnérabilité, un sous-système du noyau, ou un
        protocole réseau.
  \item \textbf{Échafaudage de code}. Générer du code de départ pour
        un harnais, un fuzzer, ou un squelette de PoC --- puis le
        modifier et l'étendre vous-même.
  \item \textbf{Débogage}. Expliquer des messages d'erreur, la sortie
        d'un désassembleur, ou un comportement de crash inattendu.
  \item \textbf{Résumé de documentation}. Condenser une longue RFC ou
        un article technique en notes exploitables.
  \item \textbf{Traduction}. Traduire du contenu technique entre
        l'anglais et le français (les deux langues sont utilisées
        dans ce cursus).
  \item \textbf{Rédaction de rapports}. Demander de l'aide pour
        structurer un rapport de bug bounty --- mais remplir le
        contenu technique vous-même.
\end{itemize}

\section{Utilisations Interdites}

\begin{itemize}
  \item \textbf{Génération complète d'exploit}. Demander à l'IA
        d'écrire l'exploit complet d'un exercice, puis le soumettre
        sans le comprendre.
  \item \textbf{Contournement des contraintes de l'exercice}. Utiliser
        l'IA pour contourner une restriction que l'exercice impose
        explicitement (p.\ ex.\ «~n'utilisez pas de scanners
        automatisés~»).
  \item \textbf{Découverte automatisée de vulnérabilités}. Exécuter des
        outils de scan pilotés par IA qui remplacent l'analyse manuelle
        que vous êtes censé pratiquer.
  \item \textbf{Rédaction des grilles d'évaluation}. Vous n'êtes pas
        autorisé à utiliser l'IA pour générer des critères d'évaluation
        par les pairs ou des réponses.
  \item \textbf{Plagiat sans attribution}. Copier le résultat de l'IA
        sans le citer (voir Section~\ref{sec:ai-citation}).
\end{itemize}

\section{Responsabilité Personnelle}

Lorsque vous demandez à l'IA de générer du code d'exploit, vous
devenez responsable de chaque octet de ce résultat avant de
l'utiliser contre une cible.

\begin{enumerate}
  \item \textbf{Lisez le code généré} ligne par ligne.
  \item \textbf{Identifiez toutes les hypothèses} faites par l'IA
        (architecture, version du noyau, drapeaux du compilateur,
        topologie réseau).
  \item \textbf{Testez en isolation} avant d'exécuter contre une
        cible d'exercice.
  \item \textbf{Comprenez les modes d'échec}. Si l'exploit généré
        plante ou se bloque, vous devez diagnostiquer et corriger
        le problème vous-même.
\end{enumerate}

\hint{Si l'IA génère du code que vous ne comprenez pas, demandez-lui
de l'expliquer. Si vous ne comprenez toujours pas après
l'explication, le code n'est pas prêt à être utilisé. Passez à
une approche plus simple que vous maîtrisez entièrement.}

\section{Citation Obligatoire de l'IA}
\label{sec:ai-citation}

Chaque fichier contenant du code généré ou assisté par l'IA doit
inclure un commentaire de citation. Le format est~:

\begin{minted}{python}
# Generated with assistance from <tool> on <date>.
# Prompt: <one-line summary of the prompt>
# Verified by: <method used to verify the output>
\end{minted}

\subsection*{Exemples}

\textbf{Python}~:
\begin{minted}{python}
# Generated with assistance from Claude 3.5 Sonnet on 2025-03-14.
# Prompt: "Write a Python script that sends HTTP requests with
# modified price fields to /api/checkout/pay"
# Verified by: manual review of each request and response
\end{minted}

\textbf{C}~:
\begin{minted}{c}
/* Generated with assistance from GitHub Copilot on 2025-03-14.
 * Prompt: "Implement a stack-based buffer overflow with
 * ret2libc for x86_64 Linux"
 * Verified by: compiled with -fno-stack-protector and
 * tested against local binary with ASLR disabled */
\end{minted}

\textbf{Markdown (rapport)}~:
\begin{minted}{html}
<!--
Generated with assistance from ChatGPT-4 on 2025-03-14.
Prompt: "Structure a bug bounty report for IDOR with
sections: summary, steps, impact, mitigation"
Verified by: all technical content written manually,
AI only suggested structure
-->
\end{minted}

\subsection*{Interactions multiples avec l'IA}

Si vous avez utilisé plusieurs outils d'IA, ou itéré à travers
plusieurs invites, citez chaque interaction significative. Un
groupe d'invites connexes peut être résumé~:

\begin{minted}{html}
<!--
AI assistance used in this file:
- Claude 3.5 Sonnet (2025-03-14): harness skeleton generation
- ChatGPT-4 (2025-03-14): crash triage help
- Copilot (2025-03-15): inline autocomplete for loop logic
All generated code reviewed and modified manually.
-->
\end{minted}

\section{Évaluation par les Pairs et l'IA}

Lors de l'évaluation par les pairs, l'évaluateur peut~:

\begin{itemize}
  \item Vous demander d'\textbf{expliquer tout code généré par l'IA}
        ligne par ligne.
  \item Vous demander de \textbf{reproduire le raisonnement} sans
        l'aide de l'IA.
  \item Demander à voir votre \textbf{historique d'invites} pour un
        exercice donné.
  \item Déduire des points si la citation de l'IA est manquante ou
        si vous ne pouvez pas expliquer le code que l'IA a généré.
\end{itemize}

\begin{warning}
  Présenter le code généré par l'IA comme votre propre travail
  (défaut de citation) est du plagiat. Si l'évaluateur soupçonne
  du plagiat, il peut saisir l'équipe pédagogique pour examen.
\end{warning}

\section{Rappel Éthique}

\begin{itemize}
  \item Les modèles d'IA peuvent produire du code incorrect,
        dangereux, ou illégal. Vous êtes \textbf{responsable} de
        vérifier chaque résultat avant exécution.
  \item Les modèles d'IA entraînés sur la recherche publique en
        sécurité peuvent reproduire textuellement du code provenant
        de projets sous licence GPL. Vérifiez les licences avant
        d'intégrer ces éléments dans votre travail.
  \item Ne saisissez pas de code sensible ou propriétaire dans des
        outils d'IA hébergés dans le cloud. Utilisez des modèles
        locaux (Llama, Mistral, DeepSeek) pour les travaux impliquant
        des cibles non publiques.
  \item L'éthique de la recherche en sécurité que vous apprenez ici
        s'applique à l'utilisation de l'IA~: ne demandez pas à l'IA
        de générer des malwares, des rançongiciels, ou du code destiné
        à nuire à des systèmes en dehors du cadre de l'exercice.
\end{itemize}

\info{La maîtrise des outils d'IA est une compétence recherchée dans
la chasse aux bogues moderne. Utilisez-les sagement, citez-les
honnêtement, et conservez toujours votre propre compréhension de
chaque exploit que vous déployez.}


\chapter{Avant-propos}

Le papier ne parle pas seulement du noyau Linux. Il utilise trois
\emph{warm-ups} pour prouver que son modèle survit au changement de domaine :

\begin{center}
\begin{tabularx}{\linewidth}{|p{0.2\linewidth}|p{0.2\linewidth}|X|}
\hline
\textbf{Domaine} & \textbf{Famille} & \textbf{Pourquoi ce cas compte} \\ \hline
Node.js & Predictive Oracle & Un prédicat de chemin croit modéliser Windows ;
le système d'exploitation en sait plus que lui. \\ \hline
Apache MINA & Logical Engine & L'ordre entre initialisation et vérification
est une transition d'état, pas un simple \texttt{if} manquant. \\ \hline
FreeBSD / XNU & Plusieurs & Le papier revendique une portée cross-kernel ;
il faut vérifier ce qui est réellement documenté. \\ \hline
\end{tabularx}
\end{center}

\begin{info}
Un \emph{warm-up} n'est pas un exercice facile. C'est un cas réduit qui
permet de voir la structure avant d'ajouter la complexité du noyau. Si vous
savez expliquer Node.js et MINA, vous devez pouvoir expliquer nf\_tables,
eBPF et io\_uring avec le même vocabulaire.
\end{info}

\chapter{Aperçu du module}

\begin{center}
\begin{tabularx}{\linewidth}{|c|l|c|c|}
\hline
\textbf{Ex} & \textbf{Titre} & \textbf{Difficulté} & \textbf{Points} \\ \hline
00 & Node.js : le prédicat et le système & Moyen & 4 \\
01 & MINA : l'ordre des opérations & Moyen & 4 \\
02 & MINA : la matrice des branches & Difficile & 4 \\
03 & FreeBSD/XNU : ce que le papier ne prouve pas & Difficile & 4 \\
04 & Synthèse : tous les exemples du papier & Difficile & 4 \\ \hline
& \textbf{Total} & & \textbf{20} \\ \hline
\end{tabularx}
\end{center}

\startexercices

\chapter{Exercise 00}
\extitle{Node.js : le prédicat et le système}
\turnindir{ex00}
\exnumber{00}
\exfiles{node-paths.md path-model.js}
\exauthorize{Node.js, un éditeur de texte, documentation Windows}
\makeheaderfiles

\texttt{path.normalize()} et \texttt{path.join()} analysent une chaîne avant
que le système de fichiers ne l'utilise. C'est un Predictive Oracle : le
prédicat décide qu'un chemin est sûr, puis le système d'exploitation lui
donne un sens différent.

Le papier décrit trois générations :

\begin{center}
\begin{tabularx}{\linewidth}{|p{0.25\linewidth}|X|X|}
\hline
\textbf{CVE / cas} & \textbf{Hypothèse du correctif} & \textbf{Angle mort} \\ \hline
CVE-2025-23084 & Un chemin qui ne commence pas par une lettre de lecteur
est relatif. & \texttt{C:..\textbackslash{}sensitive} est relatif au
répertoire courant du lecteur C:. \\ \hline
CVE-2025-27210 & Filtrer les noms de lecteur suffit à bloquer les chemins
Windows spéciaux. & \texttt{CON}, \texttt{PRN}, \texttt{AUX}, \texttt{NUL},
\texttt{COM1-9}, \texttt{LPT1-9} sont des périphériques, pas des
répertoires. \\ \hline
UNC non assigné & Filtrer les noms de périphériques suffit. & Les chemins
UNC \texttt{\textbackslash\textbackslash{}serveur\textbackslash{}partage}
restent valides et peuvent sortir du sous-arbre prévu. \\ \hline
\end{tabularx}
\end{center}

\subsection*{Mission}

Dans \texttt{node-paths.md} :

\begin{enumerate}
  \item expliquez la différence entre \texttt{path.normalize},
        \texttt{path.join} et la résolution finale par Windows ;
  \item remplissez un tableau \emph{entrée}, \emph{sortie du prédicat},
        \emph{interprétation Windows}, \emph{hypothèse violée} pour les
        trois cas ;
  \item expliquez pourquoi chaque correctif est une extension du modèle, pas
        une fermeture de la classe ;
  \item proposez un corpus de tests qui couvre lecteurs, noms réservés,
        UNC, chemins relatifs et séparateurs mélangés ;
  \item classez la chaîne : \emph{architectural scatter} ou
        \emph{incomplete-fix variant} ? Justifiez.
\end{enumerate}

Dans \texttt{path-model.js}, écrivez un petit script qui appelle uniquement
\texttt{path.win32.normalize} et \texttt{path.win32.join} sur des chaînes
fournies. Le script ne doit jamais appeler \texttt{fs}, \texttt{child\_process}
ou ouvrir un fichier.

\begin{warn}
Ne testez pas ces chemins sur un vrai disque Windows. Le but est de comparer
le modèle du prédicat au modèle du système, pas de provoquer une traversée
de répertoire.
\end{warn}

\subsection*{Sources}
\begin{itemize}
  \item NVD, CVE-2025-23084 :
        \url{https://nvd.nist.gov/vuln/detail/CVE-2025-23084}
  \item NVD, CVE-2025-27210 :
        \url{https://nvd.nist.gov/vuln/detail/CVE-2025-27210}
  \item Node.js, security releases de janvier 2025 :
        \url{https://nodejs.org/en/blog/vulnerability/january-2025-security-releases}
  \item Node.js, security releases de juillet 2025 :
        \url{https://nodejs.org/en/blog/vulnerability/july-2025-security-releases}
  \item Node.js, documentation \texttt{path} :
        \url{https://nodejs.org/api/path.html}
  \item Microsoft, nommage des fichiers et chemins Windows :
        \url{https://learn.microsoft.com/en-us/windows/win32/fileio/naming-a-file}
\end{itemize}

\chapter{Exercise 01}
\extitle{MINA : l'ordre des opérations}
\turnindir{ex01}
\exnumber{01}
\exfiles{mina-fsm.md}
\exauthorize{un éditeur de texte, sources Java, avis Apache}
\makeheaderfiles

Apache MINA utilise la désérialisation Java pour reconstruire des objets
reçus du réseau. Le papier modélise ce chemin comme une machine à états :

\begin{console}
allocate -> initialize -> resolve -> check
\end{console}

Le bug n'est pas « il manque un filtre ». Le bug est que la vérification
arrive après une transition qui a déjà exécuté du code.

\begin{center}
\begin{tabularx}{\linewidth}{|p{0.2\linewidth}|X|X|}
\hline
\textbf{CVE} & \textbf{Hypothèse} & \textbf{Transition asymétrique} \\ \hline
CVE-2024-52046 & Aucun filtre n'est nécessaire si l'application fait
confiance au flux. & Instanciation arbitraire avant tout contrôle. \\ \hline
CVE-2026-41409 & Une allowlist vérifiée dans
\texttt{AbstractIoBuffer.getObject()} protège l'instanciation. & Les
initialiseurs statiques s'exécutent avant l'évaluation de la garde. \\ \hline
CVE-2026-42778 & Le correctif de 2.0.x est complet pour toutes les
branches. & 2.1.x et 2.2.x reçoivent une seconde correction. \\ \hline
\end{tabularx}
\end{center}

\subsection*{Mission}

Dans \texttt{mina-fsm.md} :

\begin{enumerate}
  \item dessinez la machine à états \emph{allocate}, \emph{initialize},
        \emph{resolve}, \emph{check} ;
  \item placez chaque CVE sur la transition qui est mal ordonnée ;
  \item expliquez pourquoi une allowlist de classes ne suffit pas si elle
        est consultée trop tard ;
  \item identifiez l'invariant qui devrait être vrai avant toute
        initialisation ;
  \item proposez un test de régression conceptuel, sans payload.
\end{enumerate}

\begin{warn}
Ne construisez pas de flux Java sérialisé malveillant. Lisez le code, les
avis et les diffs ; décrivez la transition, pas la charge utile.
\end{warn}

\subsection*{Sources}
\begin{itemize}
  \item NVD, CVE-2024-52046 :
        \url{https://nvd.nist.gov/vuln/detail/CVE-2024-52046}
  \item NVD, CVE-2026-41409 :
        \url{https://nvd.nist.gov/vuln/detail/CVE-2026-41409}
  \item NVD, CVE-2026-42778 :
        \url{https://nvd.nist.gov/vuln/detail/CVE-2026-42778}
  \item Apache MINA, annonce CVE-2024-52046 :
        \url{https://lists.apache.org/thread/4wxktgjpggdbto15d515wdctohb0qmv8}
  \item Apache MINA, annonce CVE-2026-41409 :
        \url{https://lists.apache.org/thread/9ddvsq6c4l5bhwq8l14sob4f8qjvx5c9}
  \item Apache MINA, annonce CVE-2026-42778 :
        \url{https://lists.apache.org/thread/fhlx5k91hrkgyzh7yk1nghrn3k27gxy0}
  \item Documentation Java sur les filtres de désérialisation :
        \url{https://docs.oracle.com/en/java/javase/21/core/serialization-filtering1.html}
\end{itemize}

\chapter{Exercise 02}
\extitle{MINA : la matrice des branches}
\turnindir{ex02}
\exnumber{02}
\exfiles{branch-matrix.md}
\exauthorize{un éditeur de texte, avis Apache, historique des versions}
\makeheaderfiles

Le cas MINA est un \emph{Logical Engine} réparti sur plusieurs branches de
maintenance. Une correction peut être complète sur une branche et incomplète
sur une autre, sans nouveau chemin de code.

\subsection*{Mission}

Remplissez \texttt{branch-matrix.md} avec les branches maintenues et les
versions de correction connues :

\begin{center}
\begin{tabularx}{\linewidth}{|p{0.18\linewidth}|p{0.18\linewidth}|p{0.2\linewidth}|X|}
\hline
\textbf{Branche} & \textbf{Version corrigée} & \textbf{CVE couverte} &
\textbf{Preuve / incertitude} \\ \hline
2.0.x & & & \\ \hline
2.1.x & & & \\ \hline
2.2.x & & & \\ \hline
\end{tabularx}
\end{center}

Puis répondez :

\begin{enumerate}
  \item Pourquoi la version 2.0.28 ne prouve-t-elle pas que 2.1.11 est
        complète ?
  \item Quelle hypothèse de release le correctif de 2.0.x encode-t-il ?
  \item Comment vérifier qu'une branche a réellement reçu le déplacement de
        la garde, et pas seulement un changement de version ?
  \item Pourquoi ce cas est-il un \emph{incomplete-fix variant} et non un
        \emph{architectural scatter} ?
  \item Quelle question poser à un mainteneur avant de déclarer la classe
        fermée ?
\end{enumerate}

\subsection*{Sources}
\begin{itemize}
  \item Phrack, \emph{Security is failure: Every Patch Is a Guess},
        section 13.2.
  \item Apache MINA, avis CVE-2026-41409 et CVE-2026-42778 (mêmes liens que
        l'exercice 01).
  \item NVD, CVE-2026-41409 et CVE-2026-42778.
  \item Apache MINA, page de téléchargement et d'archives :
        \url{https://mina.apache.org/mina-project/downloads.html}
\end{itemize}

\chapter{Exercise 03}
\extitle{FreeBSD/XNU : ce que le papier ne prouve pas}
\turnindir{ex03}
\exnumber{03}
\exfiles{cross-kernel.md}
\exauthorize{un éditeur de texte, NVD, avis FreeBSD, sources XNU}
\makeheaderfiles

La section 18.11 du papier reconnaît une limite de portée : le framework a
été développé sur Linux, puis validé sur Windows. Les mappings FreeBSD et
XNU sont documentés dans l'archive, pas démontrés dans le corps de
l'article.

Deux cas FreeBSD sont cités :

\begin{itemize}
  \item CVE-2026-42511 : \texttt{dhclient} ;
  \item CVE-2026-4747 : NFS.
\end{itemize}

Le papier mentionne aussi XNU avec IOKit \texttt{IOUserClient} et Mach IPC
comme surfaces de travail, sans en faire une étude de cas complète.

\subsection*{Mission}

Dans \texttt{cross-kernel.md} :

\begin{enumerate}
  \item pour chaque cas FreeBSD, identifiez le composant, la famille
        probable et l'hypothèse que le correctif encode ;
  \item séparez explicitement \emph{ce que le papier affirme}, \emph{ce que
        l'avis public affirme} et \emph{ce que vous inférez} ;
  \item pour XNU, listez deux coutures d'audit (IOKit, Mach IPC) et la
        question de cycle de vie ou de frontière correspondante ;
  \item expliquez pourquoi une classification cross-kernel peut être un
        \emph{faux positif structurel} ;
  \item concluez : le papier démontre-t-il la portée cross-kernel, ou
        propose-t-il une hypothèse de portée ?
\end{enumerate}

\begin{warn}
N'utilisez pas ces cas pour affirmer qu'une CVE est exploitable. Le papier
lui-même rappelle que le framework produit des candidats, pas des
vulnérabilités confirmées.
\end{warn}

\subsection*{Sources}
\begin{itemize}
  \item Guide de lecture du papier :
        \texttt{00-guide/PAPER-GUIDE-paper/subject/PAPER-GUIDE-paper\_fr.pdf}.
  \item NVD, CVE-2026-42511 :
        \url{https://nvd.nist.gov/vuln/detail/CVE-2026-42511}
  \item NVD, CVE-2026-4747 :
        \url{https://nvd.nist.gov/vuln/detail/CVE-2026-4747}
  \item FreeBSD, avis de sécurité :
        \url{https://www.freebsd.org/security/advisories/}
  \item Sources XNU :
        \url{https://github.com/apple-oss-distributions/xnu}
  \item Phrack, section 18.11, \emph{What the Framework Does Not Find}.
\end{itemize}

\chapter{Exercise 04}
\extitle{Une carte exhaustive de tous les exemples du papier}
\turnindir{ex04}
\exnumber{04}
\exfiles{paper-coverage.md}
\exauthorize{un éditeur de texte, le papier, les modules, NVD}
\makeheaderfiles

Le papier contient plus d'exemples que les six CVE des études de cas
principales. Cette carte sert à vérifier qu'aucun exemple n'a été oublié et
à distinguer les preuves des extrapolations.

\subsection*{Mission}

Remplissez \texttt{paper-coverage.md} pour les CVE suivantes. Indiquez la
section du papier, la famille, le mode de mutation et le module EPIAG qui
la couvre.

\small
\renewcommand{\arraystretch}{0.95}
\begin{center}
\begin{tabularx}{\linewidth}{|p{0.18\linewidth}|p{0.1\linewidth}|p{0.18\linewidth}|X|}
\hline
\textbf{CVE} & \textbf{Section} & \textbf{Famille} & \textbf{Module / statut} \\ \hline
CVE-2023-35001 & 6.2 & Predictive Oracle & M-EPI-003, A-EPI-002 \\ \hline
CVE-2024-1086 & 6.2 & Predictive Oracle & M-EPI-003, A-EPI-002 \\ \hline
CVE-2023-4244 & 6.3 & Logical Engine & M-EPI-000, M-EPI-003 \\ \hline
CVE-2024-1085 & 6.3 & Logical Engine & M-EPI-000, M-EPI-003 \\ \hline
CVE-2025-23084 & 13.1 & Predictive Oracle & A-EPI-006 \\ \hline
CVE-2025-27210 & 13.1 & Predictive Oracle & A-EPI-006 \\ \hline
UNC non assigné & 13.1 & Predictive Oracle & A-EPI-006 \\ \hline
CVE-2024-52046 & 13.2 & Logical Engine & A-EPI-006 \\ \hline
CVE-2026-41409 & 13.2 & Logical Engine & A-EPI-006 \\ \hline
CVE-2026-42778 & 13.2 & Logical Engine & A-EPI-006 \\ \hline
CVE-2020-27194 & 14.2 & Predictive Oracle & M-EPI-004, A-EPI-003 \\ \hline
CVE-2021-3490 & 14.3 & Predictive Oracle & M-EPI-004, A-EPI-003 \\ \hline
CVE-2022-23222 & 14.4 & Predictive Oracle & M-EPI-004, A-EPI-003 \\ \hline
CVE-2023-2163 & 14.5 & Predictive Oracle & M-EPI-004, A-EPI-003 \\ \hline
CVE-2021-20226 & 15.2 & Concurrent Proxy & M-EPI-004, A-EPI-003 \\ \hline
CVE-2021-3491 & 15.2 & Concurrent Proxy & M-EPI-004, A-EPI-003 \\ \hline
CVE-2023-2598 & 15.3 & Concurrent Proxy & M-EPI-004, A-EPI-003 \\ \hline
CVE-2024-0582 & 15.3 & Concurrent Proxy & M-EPI-004, A-EPI-003 \\ \hline
CVE-2025-33073 & 16.2 & Logical Engine & A-EPI-004 \\ \hline
CVE-2025-58726 & 16.2 & Logical Engine & A-EPI-004 \\ \hline
CVE-2026-24294 & 16.2 & Logical Engine & A-EPI-004 \\ \hline
CVE-2026-26128 & 16.2 & Logical Engine & A-EPI-004 \\ \hline
CVE-2026-26119 & 16.2 & Logical Engine & A-EPI-004 \\ \hline
CVE-2026-42511 & 18.11 & Logical Engine ? & A-EPI-006 \\ \hline
CVE-2026-4747 & 18.11 & Concurrent Proxy ? & A-EPI-006 \\ \hline
\end{tabularx}
\end{center}
\normalsize

Puis ajoutez trois paragraphes :

\begin{enumerate}
  \item \textbf{Démontré} : quels cas ont une analyse détaillée dans le
        corps du papier ?
  \item \textbf{Observé} : quels cas viennent du corpus de 62 CVE ou de
        l'archive, sans analyse complète ?
  \item \textbf{Extrapolé} : quels cas servent à revendiquer une portée
        cross-kernel sans preuve détaillée ?
\end{enumerate}

Enfin, répondez à la question de soutenance : « Si vous deviez défendre ce
papier, quelle serait votre affirmation la plus forte et votre concession
la plus honnête ? »

\subsection*{Sources}
\begin{itemize}
  \item Phrack, \emph{Security is failure: Every Patch Is a Guess},
        sections 6, 13, 14, 15, 16 et 18.11.
  \item Dépôt compagnon : \texttt{validation/dataset.csv} et
        \texttt{validation/methodology.md}.
  \item NVD : \url{https://nvd.nist.gov/}
  \item Les modules EPIAG cités dans le tableau.
\end{itemize}

\chapter{Rendu et évaluation}

\begin{console}
A-EPI-006-warmups-cross-domain/
|-- ex00/
|   |-- node-paths.md
|   +-- path-model.js
|-- ex01/mina-fsm.md
|-- ex02/branch-matrix.md
|-- ex03/cross-kernel.md
+-- ex04/paper-coverage.md
\end{console}

\begin{center}
\begin{tabularx}{\linewidth}{|l|X|c|}
\hline
\textbf{Critère} & \textbf{Attendu} & \textbf{Points} \\ \hline
Node.js & Prédicat, sémantique Windows et trois angles morts & 4 \\ \hline
MINA & Ordre des transitions et invariant explicites & 4 \\ \hline
Branches & Matrice de versions et preuves sourcées & 4 \\ \hline
Cross-kernel & Distinction démontré / observé / extrapolé & 4 \\ \hline
Synthèse & Tous les exemples couverts, concession honnête & 4 \\ \hline
\end{tabularx}
\end{center}

\chapter{Sources générales}
\begin{itemize}
  \item Phrack, \emph{Security is failure: Every Patch Is a Guess},
        sections 13, 18.9 et 18.11.
  \item Node.js, avis de sécurité 2025 et documentation \texttt{path}.
  \item Apache MINA, avis CVE-2024-52046, CVE-2026-41409 et
        CVE-2026-42778.
  \item FreeBSD, avis de sécurité et NVD pour CVE-2026-42511 et
        CVE-2026-4747.
  \item Dépôt compagnon : \texttt{cross-kernel/}, \texttt{validation/} et
        \texttt{rules/}.
\end{itemize}

\end{document}
